\begin{titlepage}
\vfill
\begin{center}
{\huge
%Nachrichtentechnik\\

\vspace*{0.5cm}
\ifthenelse {\boolean{bSS}}
{Signalverarbeitung I}
{Signalverarbeitung II}
}

\vspace*{2cm}
\ifthenelse {\boolean{bSS}}
{{\large Vorlesungsscript WS 2019 / 2020 }}
{{\large Vorlesungsscript SS 2012}}

\vfill
%\includegraphics[width = 4cm]{IHA_Logo}


{\copyright Autor: Prof. Dr. Jörg Bitzer \\ \today \\ 
\ifthenelse {\boolean{bSS}}
{{für SV I Version 0.9552 \\}}{{ für SV II  Version 0.821}}}

Jade Hochschule, Fachbereich Bauwesen und Geoinformation, Abteilung Technik und Gesundheit für Menschen 

\end{center}
Fehler bzw. Änderungs-Log:
\begin{itemize}
	\item Version 0.9542 vom 29.10.2017: Fehler berichtigt Abschnitt Rekursion N ersetzt durch M und N-1 in Formel durch M
	\item Version 0.9543 vom 12.11.2017: zTrafo: nicht-kausale Folge richtig definiert: ACHTUNG: In der Bildüberschrift noch nicht geändert
	\item Version 0.9544 vom 20.11.2017: Beim Stabilitätsdreieck die Bedingung auf $a_2^2 <1$ geändert (Quadrat ergänzt).
	\item Version 0.9545 vom 26.11.2017: Begin Spektren (5.1-5.3) $\Omega$ durch $\Omega_0$ ersetzt, da eine spezifische Frequenz gemeint ist .
	\item Version 0.9546 vom 14.12.2017: Fensterfunktionen mit $4\pi$ statt $2\pi$.
	\item Version 0.9547 vom 13.11.2018: Erzeugung Dreieck cos auf sin geändert. Lösung war fehlerhaft.
	\item Version 0.9548 vom 13.11.2018: Abb 2.22 Werte zu Definitionsbereich geändert und Bsp Linearität klarer dargestellt. Es fehlt noch im oberen Teil des Bildes das xmix.
	\item Version 0.955 vom 07.12.2018: Phasenberechnung aus PN geändert.
	\item Version 0.9551 vom 05.09.2019: Girod, Rabenstein als Ideengeber bei Pol-Nullstellendarstellung 
	\item Version 0.9551 vom 05.09.2019: Fibonacci Beispiel entfernt und Stil geändert in zTrafo Anfang
\end{itemize}


Nützliche Literatur für die Module SV I und SV II und kurze Beschreibung:
\begin{itemize}
\item Girod, Rabenstein, Stenger, ''Einführung in die Systemtheorie'', Dieses Buch gibt es
in mehreren Ausgaben in der Bibliothek und ist eine anschauliche Einführung in die Systemtheorie.
Das Buch geht über den Stoff der Vorlesung hinaus, ist aber sehr anschaulich.
\item Kammeyer, Kroschel, ''Digitale Signalverarbeitung'', Eines der Standardwerke in der deutschen
SV Literatur. Gut geschrieben, Matlab Beispiele, beinhaltet theoretische Hintergründe zu einigen wichtigen Punkten der Vorlesung
\item Proakis, Manolakis, ''Digital Signal Processing'', Meine Empfehlung, ein hervorragendes Buch, Sie lernen außerdem gleich die englischen Fachbegriffe
\item Orfanidis, ''Introduction to digital signal processing'', Schönes Einsteigerbuch mit Beispielen aus dem Audiobereich
\item Martin Meyer, ''Grundlagen der Informationstechnik - Signale, Systeme und Filter'', 1. Auflage, Vieweg-Verlag,
\item Vary, Heute, Hess, ''Digitale Sprachsignalverarbeitung'', Geht inhaltlich deutlich über die Vorlesung hinaus. Ist für alle interessant, die in diesem Bereich einmal Arbeiten wollen.
\item Zölzer, ''Digitale Audiosignalverarbeitung'', Spezialbuch für alle Interessierten.
\end{itemize}


\vfill
\end{titlepage}
